% Transcription — fonds Alexandre Grothendieck, shelfmark 73, batch 2 (pages 21-28).
% Established from the facsimile archives/batches/73/batch-02.pdf, prepared
% from a copy of 73.pdf fetched through the project's own relay
% (https://grothendieck-relay.onrender.com/source/73.pdf), not uploaded by the
% user; per the skill transcribe-grothendieck. The batch file carries no cover
% sheet: PDF page k is archivists' page 20+k. The folder is twenty-eight pages;
% this is the second and last batch. Batch 1 (pages 1-20) is being transcribed
% in parallel by another pass and did not exist when this one was made, so the
% notation below was fixed from these eight pages alone.
%
% Pass: Opus 5.5 (claude-opus-5-5), 2026-09-26 — first pass, unchecked against the pages by a human.
%
% Numbering. The pencilled numbers bottom-left of these leaves read 28 to 35,
% not 21 to 28; the last leaf of batch 1 carries a pencilled 27 and its first
% a pencilled 3, so that pencilled series is not continuous with the PDF's
% count and is not the one the catalogue's « 28 pages » counts. \page follows
% the PDF count (20+k), as for every other folder; the discrepancy is recorded
% here only.
%
% Dating and title. The dating is the folder's own, copied verbatim from the
% catalogue; the group « Géométrie et topologie combinatoire » (69-89,
% 1976-[vers 1986]) holds it, and since the folder has a date of its own the
% group's range is not added. The title « Polygones : notes manuscrites
% (s.d.), lettre (s.d.). » is the catalogue's; nothing in it is bracketed.
%
% Pages skipped: none. Pages summarised: none. There is no letter in this
% batch, by him or by anyone else: all eight leaves are his working notes, in
% blue-black ink, one hand throughout. The letter the catalogue mentions must
% lie in batch 1.
%
% His pagination: every leaf carries his own page number top right, mostly
% circled: 9 (p. 21), 10 (p. 22), 11 (p. 23), 12 (p. 24), 13 (p. 25),
% 14 (p. 26), 15 (p. 27), 16 (p. 28). The run therefore begins in batch 1
% (whose last leaf, p. 20, carries his « 8 ») and is recorded once per page
% with a \note.
%
% What the batch is about. A closed polygon in a plane V over a field k is
% given by its side vectors u_i, i in Z/nZ, and each vertex carries a linear
% relation between four consecutive sides,
%    lambda_i u_{i-1} - lambda_{i+1}(u_i + u_{i+1}) + lambda_{i+2} u_{i+2} = 0.
%   21-22  Solving the recurrence: u_i = lambda_{i-1}/(lambda_3...lambda_i)
%          (P_i u_0 + Q_i u_1 + R_i u_2), the three-term recurrence for P, Q,
%          R (boxed), P_3 = -lambda_1, Q_3 = R_3 = lambda_2, P_i, Q_i, R_i
%          forms of degree i-3; the closing conditions u_n = u_0, u_{n+1} =
%          u_1, u_{n+2} = u_2 give three boxed equations (1)-(3), hence a
%          system with coefficients A_i^n, B_i^n, C_i^n.
%   23     Those coefficients are integral polynomials in indeterminates
%          Lambda; NB: lambda_1, lambda_2, lambda_3 fixed determine the rest;
%          if some u_i, u_{i+1}+u_{i+2}, u_{i+3} are not collinear the
%          (lambda_i) are determined up to a scalar. « Mais quand ... ! » —
%          the degenerate case where all such triples are collinear is taken
%          up with parameters alpha, beta, gamma, delta.
%   24     The degenerate case computed: u_{i+3} = alpha_i u_i, then with
%          beta's; u_1 = u_4, hence u_i = u_{i+3} and n = 0 mod 3.
%   25     « Proposition »: for u_* with non-zero sides, N(u_*) the space of
%          admissible (lambda_i), equivalence of (a) some (lambda_i) with all
%          lambda_i non-zero and dim_k N(u_*) >= 2, (b) every u_i,
%          u_{i+1}+u_{i+2}, u_{i+3} collinear, (c) u_{i+3} = u_i and u_{i+1} +
%          u_{i+2} = -u_i, (d) u_i = u_{i+3} and u_0+u_1+u_2 = 0. Heavily
%          reworked; the statement is given as it finally reads, with the
%          cancelled layers struck.
%   26-27  « Dém. »: for k infinite, N(u_*) -> k^3, (lambda_i) -> (lambda_1,
%          lambda_2, lambda_3) is injective, 1 <= dim N(u_*) <= 3, dim 3 is
%          excluded, dim N(u_*) <= 3 - rg(u_0, u_1+u_2, u_3); a) => b),
%          b) <=> c) <=> d), and the 3-periodic (lambda_i) form a space of
%          dimension 2. « Dém. d) => a) » is announced and not written.
%   28     « Théorème »: the trichotomy for N(u_*), labelled in the margin
%          « Cas général » (a) dim N(u_*) = 0, « Cas spécial I » (b_1) dim 1,
%          « Cas hyperspécial » (b_2) dim 2 (u_i = u_{i+3}, u_0+u_1+u_2 = 0),
%          « Cas spécial II » (c) N(u_*) inside a coordinate hyperplane; with a
%          criterion for (c) and a diagonal marginal note that is largely lost.
%
% Notation, fixed once for the batch. Double-struck letters as \mathbf
% (\mathbf{Z}, as in folders 75 and 86); his capital Lambda for the
% indeterminates of the polynomial rings kept as \Lambda; his « u_* » for the
% sequence of sides kept with an asterisk subscript; « // » for « parallel
% to » as \parallel; « rg » as \operatorname{rg}. He writes a product or a
% list with a long dash between its ends (« lambda_3 -- lambda_{n-1} »,
% « lambda_1, -- lambda_{n-2} »); the product is set \lambda_3\cdots
% \lambda_{n-1} and the argument list \lambda_1,\ldots,\lambda_{n-2}, with
% one note at the first occurrence. His circled item letters (a)-(d), b_1,
% b_2 are given as (a) ... with a note. His boxed formulas are given in
% display with a note saying they are boxed.
%
% Hardest pages: 25 (the Proposition, rewritten three times in place, whose
% cancelled layers are mostly illegible), 26 (the long parenthesis on the
% kernel of N(u_*) -> k^3 and its translates in the union of the hyperplanes
% H_i) and 28 (the bottom third and the diagonal marginal note). Readings to
% check first: the superscripts n on A, B, C (p. 22-23); « à la
% multiplication près » and the closing sentence « Mais quand ... ! » on
% p. 23; the indices of u_7, u_8, u_9 on p. 24 (they do not follow the rule
% u_{i+3} = alpha_i u_i written above them, and are kept as read); the
% object phi is defined on at p. 28.
\documentclass[11pt,a4paper]{article}
% Le préambule commun à toutes les transcriptions du fonds Grothendieck.
%
% Il tient en une idée : **l'incertitude se compose, elle ne se résout pas.**
% Une transcription qui lisse un mot douteux a détruit la seule chose pour
% laquelle on la faisait — un lecteur doit pouvoir savoir, à chaque endroit,
% ce qui était sur le feuillet et ce qui a été ajouté. Les six macros qui
% suivent sont donc l'appareil critique, et non de la décoration.
%
% Les mêmes commandes sont reconnues par `scripts/render.mjs`, qui produit la
% vue de lecture du site. Une macro ajoutée ici doit l'être aussi là-bas,
% faute de quoi le rendu échouera bruyamment — ce qui est voulu : un rendu
% silencieusement faux est pire qu'un rendu absent.

\ProvidesPackage{grothendieck}[2026/08/08 Transcriptions du fonds Grothendieck]

\usepackage{amsmath,amssymb,amsthm}
\usepackage{mathrsfs}
\usepackage{stmaryrd}
% Les diagrammes commutatifs sont partout dans ce fonds : c'est la langue
% même de la géométrie algébrique de Grothendieck, et une transcription qui
% les rendrait en prose ne transcrirait rien.
\usepackage{tikz-cd}
% Français par défaut — c'est la langue des feuillets — mais l'anglais est
% chargé aussi : la traduction et le résumé sont des documents anglais, et la
% typographie française (espaces avant les deux-points) y serait une faute.
% Ces documents basculent par \selectlanguage{english} en tête de corps.
\usepackage[english,french]{babel}
\usepackage{fontspec}
\usepackage{geometry}
\geometry{a4paper,margin=25mm,marginparwidth=28mm}
\usepackage{marginnote}
\usepackage{xcolor}
% ulem, not soul. `soul`'s \st and \ul are built on a character-by-character
% scan that XeLaTeX's font machinery breaks: the document fails with an
% undefined control sequence rather than degrading. ulem does the same two jobs
% and is Unicode-engine safe. `normalem` stops it hijacking \emph, which the
% transcriptions use for Grothendieck's own emphasis.
\usepackage[normalem]{ulem}
\usepackage{hyperref}
\hypersetup{colorlinks=true,linkcolor=black,urlcolor=black,pdfborder={0 0 0}}

% --- Métadonnées du lot -----------------------------------------------------
% Reprises telles quelles par le script de rendu, qui les affiche en tête de
% la vue de lecture. Elles fixent ce que le fichier prétend couvrir : sans
% elles, une transcription flotte sans point d'attache dans un fonds de seize
% mille feuillets.

\newcommand{\folder}[1]{\gdef\grothFolder{#1}}
\newcommand{\batch}[1]{\gdef\grothBatch{#1}}
\newcommand{\leaves}[2]{\gdef\grothFirst{#1}\gdef\grothLast{#2}}
\newcommand{\foldertitle}[1]{\gdef\grothFoldertitle{#1}}
\gdef\grothFolder{?}\gdef\grothBatch{?}\gdef\grothFirst{?}\gdef\grothLast{?}\gdef\grothFoldertitle{}

% --- Le résumé --------------------------------------------------------------
%
% Il ouvre la lecture modernisée, sous le titre « Résumé », et il est composé
% en petit. Ce n'est pas un ornement : c'est le seul passage du document qui
% ne soit pas des mathématiques, et un lecteur doit pouvoir le reconnaître
% avant de l'avoir lu — puis le sauter s'il connaît déjà le sujet, ou s'y
% arrêter s'il ne le connaît pas. Le filet qui le suit marque la reprise du
% corps.

\newenvironment{resume}
  {\small\setlength{\parskip}{0.45em}}
  {\par\medskip\hrule\medskip}

% --- Mots-clés ---------------------------------------------------------------
%
% En anglais, en clôture du résumé de la lecture modernisée : le vocabulaire
% moderne sous lequel chercher ce que les feuillets construisent. Le manifeste
% du site (`npm run manifest`) les extrait pour étiqueter la cote entière —
% cette ligne est la source unique des tags, il n'y a pas de fichier à côté.
\newcommand{\keywords}[1]{\par\smallskip\noindent
  {\footnotesize\textsc{Keywords} --- \textit{#1}}\par}

% --- L'appareil critique ----------------------------------------------------

% Le numéro de feuillet, en marge. C'est l'ancre qui permet de régler un
% désaccord sur une formule en regardant *un* feuillet plutôt qu'une chemise
% de six cents. Le site s'en sert aussi pour tourner les pages du fac-similé
% quand on fait défiler la transcription.
\newcommand{\leaf}[1]{%
  \marginnote{\footnotesize\color{gray!70}\textsf{#1}}%
  \label{leaf:#1}\ignorespaces}

% Illisible : on ne devine pas. Un mot inventé qui se lit comme les autres est
% le pire résultat possible.
\newcommand{\ill}{\textcolor{red!60!black}{[\,\ldots\,]}}

% Lecture douteuse : le mot est donné, et signalé comme incertain.
\newcommand{\uncertain}[1]{\uline{#1}}

% Ajout de l'éditeur — une lettre manquante, un mot sous-entendu.
\newcommand{\add}[1]{\textcolor{blue!50!black}{[#1]}}

% Biffé par Grothendieck lui-même. Ce qu'il a rayé fait partie du document :
% une rature dit souvent où la pensée a changé de direction.
\newcommand{\struck}[1]{\sout{#1}}

% Note du transcripteur, sur l'état matériel du feuillet.
\newcommand{\note}[1]{\par\smallskip{\small\color{gray!60!black}\textit{[#1]}}\par\smallskip}

% Note marginale de Grothendieck, distincte du corps du texte.
\newcommand{\marginal}[1]{\par\smallskip\noindent\hspace{1em}%
  {\small\color{gray!60!black}#1}\par\smallskip}

% --- Titre ------------------------------------------------------------------

\AtBeginDocument{%
  \begin{center}
    {\large\bfseries Fonds Alexandre Grothendieck --- cote n\textsuperscript{o}~\grothFolder}\\[2pt]
    {\itshape\grothFoldertitle}\\[4pt]
    {\small feuillets \grothFirst--\grothLast\ (lot \grothBatch)}\\[2pt]
    {\footnotesize\color{gray!60!black}Transcription établie d'après le fac-similé numérisé
     par l'Université de Montpellier.\\
     Les lectures incertaines sont soulignées, les illisibles notées [\,\ldots\,],
     les ajouts entre crochets.}
  \end{center}
  \vspace{1em}}

\endinput


\folder{73}
\batch{2}
\pages{21}{28}
\dating{[à partir de 1976]}
\watermark{Édition de démonstration}
\foldertitle{Polygones : notes manuscrites (s.d.), lettre (s.d.).}

\begin{document}

\page{21}
\note{p. 9 de sa pagination, numéro cerclé en haut à droite. Les notations
viennent des pages précédentes (lot 1) : $u_i$ ($i \in \mathbf{Z}/n\mathbf{Z}$)
sont les côtés d'un polygone fermé, et chaque sommet donne une relation
$\lambda_{i-1} u_{i-2} - \lambda_i (u_{i-1} + u_i) + \lambda_{i+1} u_{i+1} = 0$.
Il écrit un produit ou une suite d'arguments avec un long tiret entre les
extrémités (« $\lambda_3$ -- $\lambda_{n-1}$ », « $\lambda_1$, -- $\lambda_{n-2}$ ») :
on le rend ici et dans la suite par $\lambda_3 \cdots \lambda_{n-1}$ et
$\lambda_1, \ldots, \lambda_{n-2}$}

\[
u_3 = \lambda_3^{-1} \, (P_3 u_0 + Q_3 u_1 + R_3 u_2)
\]
\note{ligne ajoutée au-dessus de la suivante ; l'exposant $-1$ est
\uncertain{lu}}
\[
u_i = \frac{\lambda_{i-1}}{\lambda_3 \lambda_4 \cdots \lambda_i}
\bigl( P_i(\lambda_1, \ldots, \lambda_{i-1})\, u_0
 + Q_i(\lambda_1, \ldots, \lambda_{i-1})\, u_1
 + R_i(\lambda_1, \ldots, \lambda_{i-1})\, u_2 \bigr)
\qquad (i \geqslant 4)
\]

\begin{align*}
\lambda_{i+1} u_{i+1} &= \lambda_i (u_i + u_{i-1}) - \lambda_{i-1} u_{i-2} \\
 &= \frac{1}{\lambda_3 \cdots \lambda_{i-2}}
    \bigl[ P_i u_0 + Q_i u_1 + R_i u_2 \bigr]
  + \frac{\lambda_i \lambda_{i-2}}{\lambda_3 \cdots \lambda_{i-1}}
    \bigl[ P_{i-1} u_0 + Q_{i-1} u_1 + R_{i-1} u_2 \bigr] \\
 &\quad - \frac{\lambda_{i-1} \lambda_{i-3}}{\lambda_3 \lambda_4 \cdots \lambda_{i-2}}
    \bigl[ P_{i-2} u_0 + Q_{i-2} u_1 + R_{i-2} u_2 \bigr]
\end{align*}
\note{un petit « 2 » est écrit au-dessus de « $R_i u_2$ » à la fin de la
deuxième ligne, sans rattachement clair}

\[
u_{i+1} = \frac{\lambda_i}{\lambda_3 \cdots \lambda_{i+1}}
\Bigl\lbrace \lambda_{i-1} \bigl[\quad\bigr]_1
 + \lambda_i \lambda_{i-2} \bigl[\quad\bigr]_2
 - \lambda_{i-1}^2 \lambda_{i-3} \bigl[\quad\bigr]_3 \Bigr\rbrace
\]

Au-dessous, une petite table :
\[
\begin{array}{ccccc}
i-3 & i-2 & i-1 & i & i+1 \\
0 & -1 & 1 & 0 & -1
\end{array}
\]

\[
\left\lbrace
\begin{array}{l}
P_{i+1} = \lambda_{i-1} P_i + \lambda_i \lambda_{i-2} P_{i-1}
 - \lambda_{i-1}^2 \lambda_{i-3} P_{i-2} \\
Q_{i+1} = \lambda_{i-1} Q_i + \lambda_i \lambda_{i-2} Q_{i-1}
 - \lambda_{i-1}^2 \lambda_{i-3} Q_{i-2} \\
R_{i+1} = \lambda_{i-1} R_i + \lambda_i \lambda_{i-2} R_{i-1}
 - \lambda_{i-1}^2 \lambda_{i-3} R_{i-2}
\end{array}
\right.
\qquad i \geqslant 6
\]
\note{les trois récurrences sont encadrées. Dans les deux premières lignes,
les lettres $P_{i-1}$, $P_{i-2}$, $Q_{i-2}$ sont écrites en surcharge d'un
$R$ (ou l'inverse). Sous le cadre, \struck{$5 \leqslant i$} est surchargé et
remplacé par $i \geqslant 6$}

\[
P_3 = -\lambda_1, \qquad Q_3 = \lambda_2, \qquad R_3 = \lambda_2
\]
\marginal{$P_i, Q_i, R_i$ \emph{formes} de degré $i-3$ ($i \geqslant 4$)}
\note{séparé des valeurs initiales par un trait oblique ; l'indice de
$\lambda_2$ dans $Q_3$ est surchargé}

\[
(\lambda_3 \cdots \lambda_{n-1}) \sum_{0}^{n-1} u_i
 = \Bigl( \sum_{i=0}^{n-1} P_i(\lambda_1, \ldots, \lambda_{n-2}) \Bigr) u_0
 + (\quad) u_1 + (\quad) u_2
\]
\note{le membre de droite est barré de traits obliques, et une accolade le
remplace par :}
\[
\widetilde{P}_{n-1}(\lambda_1, \ldots, \lambda_{n-2})\, u_0
 + \widetilde{Q}_{n-1}(\lambda_1, \ldots, \lambda_{n-2})\, u_1
 + \widetilde{R}_{n-1}(\lambda_1, \ldots, \lambda_{n-2})\, u_2
\]

\struck{$\widetilde{P}_{n-1}(\lambda_1, \ldots, \lambda_{n-2})\, u_0
 + \widetilde{Q}_{n-1}(\quad)\, u_1 + \widetilde{R}_{n-1}(\quad)\, u_2 = 0$}

\begin{align*}
\widetilde{P}_{n-1}(\lambda_1, \ldots, \lambda_{n-1})
 = (\lambda_3 \cdots \lambda_{n-1})
 &+ \bigl( \lambda_4 \lambda_5 \cdots \lambda_{n-1} P_3
 + \lambda_3 \lambda_5 \lambda_6 \cdots \lambda_{n-1} P_4 \\
 &\qquad + \lambda_4 \lambda_6 \lambda_7 \cdots \lambda_{n-1} P_5 + \cdots \bigr)
\end{align*}
\note{avant la parenthèse, un début \struck{$\sum_{i=3}^{n-1} \lambda_{i-1}
\lambda_{i+1} \cdots$} est biffé. Dans le troisième terme, le $\lambda_4$ est
écrit en surcharge d'un $\lambda_5$ \uncertain{(lecture)}}

d'où
\[
\widetilde{P}_{n-1} \in \mathbf{Z}[\Lambda_1, \ldots, \Lambda_{n-1}]
\]

\[
(0) \qquad \widetilde{P}_{n-1} u_0 + \widetilde{Q}_{n-1} u_1
 + \widetilde{R}_{n-1} u_2 = 0
\]
\note{la relation est encadrée ; le « (0) » qui la numérote est \uncertain{lu},
suivi d'un signe biffé, peut-être un $\Sigma$}

\page{22}
\note{p. 10 de sa pagination}

\struck{$u_n = \lambda P_n(\lambda_1, \ldots, \lambda_n)\, u_0 + \cdots$}
\quad ($\lambda_n = \lambda_0$)

\[
\underset{\textstyle u_0}{\underset{\shortparallel}{u_n}}
 = \frac{\lambda_{n-1}}{\lambda_3 \cdots \lambda_{n-1} \lambda_n}
 \bigl[ P_n(\lambda_1, \ldots, \lambda_{n-1}, \lambda_n)\, u_0
 + Q_n(\lambda_1, \ldots, \lambda_n)\, u_1
 + R_n(\lambda_1, \ldots, \lambda_n)\, u_2 \bigr]
\]
\note{sous chaque $\lambda_n$, un « $=$ » vertical et $\lambda_0$}

$= u_0$ i.e.
\[
(1) \qquad \bigl[ \lambda_{n-1} P_n - (\lambda_3 \cdots \lambda_{n-1} \lambda_n) \bigr] u_0
 + \lambda_{n-1} Q_n(\quad)\, u_1 + \lambda_{n-1} R_n(\quad)\, u_2 = 0
 \qquad (\lambda_n = \lambda_0)
\]
\note{encadrée. Le facteur $\lambda_{n-1}$ devant $R_n$ est ajouté en
interligne}

\[
\underset{\textstyle u_1}{\underset{\shortparallel}{u_{n+1}}}
 = \frac{\lambda_n}{\lambda_3 \cdots \lambda_{n+1}}
 \bigl[ P_{n+1}(\lambda_1, \ldots, \lambda_{n-1}, \lambda_n, \lambda_{n+1})\, u_0
 + \cdots \bigr]
 \qquad (\lambda_n = \lambda_0,\ \lambda_{n+1} = \lambda_1)
\]
$= u_1$ i.e.
\[
(2) \qquad \lambda_n P_{n+1}(\quad)\, u_0
 + \bigl( \lambda_n Q_{n+1} - (\lambda_3 \lambda_4 \cdots \lambda_{n+1}) \bigr) u_1
 + \lambda_n R_{n+1}(\quad)\, u_2 = 0
\]
\note{encadrée ; le « $= 0$ » est écrit sous $u_2$, par un « $=$ » vertical}

\[
u_{n+2} = \frac{\lambda_{n+1}}{\lambda_3 \cdots \lambda_{n+2}}
\bigl[ P_{n+2} u_0 + Q_{n+2} u_1 + R_{n+2} u_2 \bigr] = u_2
\]
\[
(3) \qquad \lambda_{n+1} P_{n+2}(\quad)\, u_0 + \lambda_{n+1} Q_{n+2}(\quad)\, u_1
 + \lambda_{n+1} R_{n+2}(-)(\lambda_3 \cdots \lambda_{n+2})\, u_2 = 0
\]
\note{encadrée, telle qu'écrite : entre $R_{n+2}$ et le produit
$(\lambda_3 \cdots \lambda_{n+2})$ il n'y a qu'un tiret entre parenthèses, là où
(1) et (2) conduiraient à « $\bigl(\lambda_{n+1} R_{n+2} - (\lambda_3 \cdots
\lambda_{n+2})\bigr) u_2$ » ; le dernier indice du produit est surchargé}

i.e. \struck{$A\,A\ A_0(\lambda_1, \ldots, \lambda_{n-2})\, u_0 + A_1($}
\[
\left\lbrace
\begin{array}{l}
A_0^n(\lambda_1, \ldots, \lambda_{n-2})\, u_0 + B_0^n(\lambda_1, \ldots, \lambda_{n-2})\, u_1
 + C_0^n(\lambda_1, \ldots, \lambda_{n-2})\, u_2 = 0 \\
A_1^n(\lambda_0, \ldots, \lambda_{n-1})\, u_0 + B_1^n(\lambda_0, \ldots, \lambda_{n-1})\, u_1
 + C_1^n(\lambda_0, \ldots, \lambda_{n-1})\, u_2 = 0 \\
A_2^n(\quad)\, u_0 + B_2^n(\quad)\, u_1 + C_2^n(\quad)\, u_2 = 0 \\
A_3^n(\quad)\, u_0 + B_3^n(\quad)\, u_1 + C_3^n(\quad)\, u_2 = 0
\end{array}
\right.
\]
\note{le système porte à gauche une marque \uncertain{$(*)$}, surchargée. Les
exposants $n$ de $A$, $B$, $C$ ne sont nets qu'à la première ligne ; ailleurs
ils sont réduits à un trait}

\page{23}
\note{p. 11 de sa pagination}

\[
\begin{array}{c}
A_i^n,\ B_i^n,\ C_i^n \in \mathbf{Z}[\Lambda_0, \ldots, \Lambda_{n-1}] \\
\cup \\
A_0^n,\ B_0^n,\ C_0^n \in \mathbf{Z}[\Lambda_1, \ldots, \Lambda_{n-2}]
\end{array}
\]
\note{le signe entre les deux lignes est un $\cup$, sans doute pour une
inclusion écrite verticalement}

\struck{Prenons $\lambda_3 = 1$, \ill{} $\lambda_i$ ($i \geqslant 4$) \ill{}
les $u_i$ \ill{} $\leqslant i \leqslant n-1$) et}

\textbf{NB} Si $\lambda_1 \lambda_2 \lambda_3$ fixés, \struck{\ill{}}
\add{$u_i$ et les} si de plus on impose donc les $\lambda_i$
($0 \leqslant i \leqslant n-1$) $\neq 0$, alors les $\lambda_i$
($4 \leqslant i \leqslant n-1$) sont déterminés de façon unique. D'ailleurs la
relation
\[
\lambda_1 u_0 - \lambda_2 (u_1 + u_2) + \lambda_3 u_3 = 0
\]
montre que, si $u_0$, $u_1 + u_2 = u'_1$, $u_3$ ne sont pas coll., les
$(\lambda_1, \lambda_2, \lambda_3)$ possibles (pour cette relation) forment un
espace de dim.~$1$, i.e.\ $(\lambda_1, \lambda_2, \lambda_3)$ est déterminé
\uncertain{à la} multiplication près par \ill{} --- donc aussi
$(\lambda_0, \ldots, \lambda_{n-1})$. Changeant la \uncertain{notation} par
permutation circulaire, on voit qu'il en est ainsi dès qu'il existe $i$, avec
$u_i$, $u_{i+1} + u_{i+2}$, $u_{i+3}$ non colin. Mais quand \ill{} \ill{}
\struck{\ill{}} \ill{} coll.\ $\forall i$ !
\note{les mots « $u_i$ et les », insérés au-dessus de la première ligne du NB,
remplacent un mot biffé ; l'endroit exact où ils se placent n'est pas sûr}

\begin{tikzcd}[column sep=small, row sep=small]
 & & s_2 \arrow[dr, "u_2"] & & \\
s_0 \arrow[r, "u_0"] & s_1 \arrow[ur, "u_1"] & & s_3 \arrow[r, "u_3"] & {}
\end{tikzcd}

\note{croquis en marge gauche du paragraphe : la ligne brisée
$s_0 s_1 s_2 s_3$, côtés $u_0, \ldots, u_3$, $u_0$ et $u_3$ horizontaux}

\marginal{OPS $u_0, u_1$ lin.\ indép.}

\[
\begin{array}{ll}
u_1 + u_2 = -\alpha_0 u_0 & \\
u_3 = \beta u_0 &
\end{array}
\qquad
\left\lbrace
\begin{array}{l}
u_2 = -\alpha u_0 - u_1 \\
u_3 = \beta u_0
\end{array}
\right.
\]
\struck{$\alpha$} $u_2 + u_3 = \gamma u_1$, \quad $u_4 = \delta u_1$ \quad i.e.
\[
\left\lbrace
\begin{array}{l}
u_3 = \alpha u_0 + (1 + \gamma) u_1 \\
u_4 = \delta u_1
\end{array}
\right.
\qquad
\begin{array}{l}
\gamma = -1 \\
\beta = \alpha
\end{array}
\]
\note{dans « $\beta = \alpha$ », l'$\alpha$ est écrit en surcharge}

\page{24}
\note{p. 12 de sa pagination}

\[
u_2 = \alpha u_0 - u_1, \qquad u_3 = \alpha u_0, \qquad u_4 = \delta u_1
\]
\note{dans les deux premières, un signe $-$ devant $\alpha$ est surchargé ou
biffé ; un point au-dessus de $u_1$}

\[
u_{i+3} = \alpha_i u_i
\]
\note{à gauche, un croquis : une ligne brisée de trois segments, sans lettres}

\[
u_0 \ / \ u_1 \ /
\]
\[
\begin{array}{l}
u_2 = -\alpha_0 u_0 - u_1 \\
u_3 = \alpha_0 u_0 \\
u_4 = \alpha_1 u_1 \\
u_5 = \alpha_2 (\alpha_0 u_0 - u_1) \\
u_6 = \alpha_3 \alpha_0 u_0 \\
u_7 = \alpha_4 \alpha_2 (-\alpha_0 u_0 - u_1) \\
u_8 = \alpha_5 \alpha_3 \alpha_0 u_0 \\
u_9 = \alpha_6 \alpha_4 \alpha_2 (-\alpha_0 u_0 - u_1)
\end{array}
\]
\note{tel qu'écrit. Pour $u_5$ la parenthèse porte un trait qui peut être un
signe $-$ devant $\alpha_0$ ; pour $u_8$ un trait relie $\alpha_5$ à $\alpha_3$.
Les indices de $u_7$, $u_8$, $u_9$ ne suivent pas la règle
$u_{i+3} = \alpha_i u_i$ écrite au-dessus ; ils ne sont pas corrigés. Un « $=$ »
isolé se trouve à gauche de la colonne}

\[
\left.
\begin{array}{l} u_0 \\ u_1 \end{array}
\right\rbrace \text{ lin.\ indép.}
\]
$\Longrightarrow (u_1, u_2)$ lin.\ indép., donc \uncertain{tous} $(u_i, u_{i+1})$
lin.\ indép.) ($\beta_0 = \alpha_0$ et
\[
\beta_1 = \alpha_1, \qquad \beta_2 = \alpha_2, \qquad \beta_3 = \alpha_3 \alpha_0
\]

\[
\begin{array}{l}
u_2 = -\beta_0 u_0 - u_1 \\
u_3 = \beta_0 u_0 \\
u_4 = \beta_1 u_1 \\
u_5 = \beta_2 u_2 \\
u_6 = \beta_3 u_0 \\
u_7 = \beta_4 u_2
\end{array}
\]
\note{dans les deux premières lignes, $\beta_0$ est écrit en surcharge d'un
$\alpha_0$}

$\underline{u_0}$, $\underline{u_2}$ lin.\ indép.
\struck{$u_1$} $= -u_2 - \beta_0 u_0$, \qquad $u_3 = \beta$
\note{ces deux lignes sont entourées d'un trait qui les relie à « $u_0$, $u_2$
lin.\ indép. »}

\[
u_3 + u_4 = \beta_0 u_0 + \beta_1 u_1 \parallel u_2 = -\beta_0 u_0 - u_1
\qquad \beta_1 = 1 \quad \text{i.e.} \quad u_1 = u_4
\]

\struck{donc} $u_i = u_{i+3}$ \qquad \struck{$u_0 u_1 u_2$}

donc $n \equiv 0 \ (3)$ \quad (\ill{} $u_i = u_{i+1} = \cdots$ $\forall i$
\uncertain{exclu}\ldots)
\[
\begin{array}{cccccccccc}
u_0 & u_1 & u_2 & u_0 & u_1 & u_2 & u_0 & u_1 & u_2 & \ldots \\
 & & & \shortparallel & \shortparallel & \shortparallel & \shortparallel & \shortparallel & \shortparallel & \\
 & & & u_3 & u_4 & u_5 & u_6 & u_7 & u_8 &
\end{array}
\]
\note{la ligne inférieure, très serrée, est \uncertain{lue} ; deux indices y
sont surchargés}

\page{25}
\note{p. 13 de sa pagination, numéro cerclé}

\[
\begin{array}{l}
u_1 + u_2 \parallel u_0 \\
u_2 + u_0 \parallel u_1 \\
u_0 + u_1 \parallel u_2
\end{array}
\qquad
\begin{array}{l}
u_2 = \alpha_0 u_0 - u_1 \\
u_2 + u_0 = (1 + \alpha_0) u_0 - u_1 \parallel u_1
 \quad \text{i.e.} \quad 1 + \alpha_0 = 0, \quad \alpha_0 = -1 \\
u_2 = -u_0 - u_1
\end{array}
\]
\note{l'$\alpha_0$ de la première ligne de droite est écrit en surcharge ; dans
« $\alpha_0 = -1$ » l'indice et le $1$ sont surchargés}

\emph{Proposition} Soit $u_* = (u_i)_{i \in \mathbf{Z}/n\mathbf{Z}}$ une
suite de $n$ vecteurs \add{\ill{}} \emph{non nuls} d'un plan vectoriel $V$,
\struck{\ill{}} \struck{(\ill{})}, et \struck{supposons qu'il existe}
\add{considérons l'ens.\ $N(u_*)$ des} $(\lambda_i)_{i \in \mathbf{Z}/n\mathbf{Z}}$
dans $k^{\mathbf{Z}/n\mathbf{Z}}$ tels que \uncertain{la suite satisfait} les
relations
\[
(*) \qquad \lambda_i u_{i-1} - \lambda_{i+1} (u_i + u_{i+1}) + \lambda_{i+2} u_{i+2} = 0
\qquad \forall i \in \mathbf{Z}/n\mathbf{Z}
\]
\note{dans la définition, « $\mathbf{Z}/n\mathbf{Z}$ » remplace un premier
indice biffé ; « considérons l'ens.\ $N(u_*)$ des » est écrit au-dessus de
« supposons qu'il existe », biffé. La marque $(*)$ est \uncertain{lue}}

\struck{Conditions équivalentes : \ill{} \ill{} \ill{} un \uncertain{vecteur}
\ill{} des}

\struck{Supposons} (a) \struck{que} $\exists\, (\lambda_i) \in N(u_*)$ tels que
\struck{\ill{} $\dim N(u_*) \geqslant 1$)} \emph{tous} les $\lambda_i$ soient
$\neq 0$, \struck{\ill{} que} et \struck{$\dim N(u_*) \geqslant 2$, il \ill{}
et \ill{} que \ill{} soit \ill{} colinéaires \ill{}} $\dim_k N(u_*) \geqslant 2$
\note{le passage biffé sur $\dim N(u_*) \geqslant 2$ est enfermé dans un
cadre barré de traits en zigzag ; au pied du cadre, « a) » suivi d'une lettre
\uncertain{$k$}. La condition (a) se lit, après corrections : il existe
$(\lambda_i) \in N(u_*)$ avec tous les $\lambda_i \neq 0$, et
$\dim_k N(u_*) \geqslant 2$. Les lettres (a) à (d) sont cerclées}

(b) $\forall i$, $u_i$, $u_{i+1} + u_{i+2}$, $u_{i+3}$ sont col.

(c) $\forall i$, \ldots
\[
\left\lbrace
\begin{array}{l}
u_{i+3} = u_i \\
u_{i+1} + u_{i+2} = -u_i \qquad \forall i
\end{array}
\right.
\]
\note{après $u_{i+3} = u_i$, une suite \struck{$u_0 + u_1 + u_2 = 0$} et un
début illisible sont biffés}

(d)
\[
\left\lbrace
\begin{array}{l}
\forall i \quad u_i = u_{i+3} \\
u_0 + u_1 + u_2 = 0
\end{array}
\right.
\]
\note{après $u_i = u_{i+3}$, une parenthèse biffée, \struck{\ill{} (3)}, sans
doute « $n \equiv 0$ (3) »}

\page{26}
\note{p. 14 de sa pagination, numéro cerclé}

\marginal{OPS $k$ infini}

\emph{Dém.} Un $(\lambda_i) \in N(u_*)$ tel que \struck{les} \uncertain{comp.}
\ill{} $\neq 0$ \uncertain{est connu} \ill{} \uncertain{quand} on connaît trois
comp.\ consécutives, disons $\lambda_1, \lambda_2, \lambda_3$ (en utilisant le
fait que les $u_i$ sont $\neq 0$), \struck{et on en conclut} i.e.
\[
\begin{array}{ccc}
N(u_*) & \longrightarrow & k^3 \\
(\lambda_i) & \longmapsto & (\lambda_1, \lambda_2, \lambda_3)
\end{array}
\]
a un noyau \uncertain{dont des} translatés \struck{\ill{}} \add{dans $N(u_*)$}
\struck{contenus dans} la réunion des hyperplans \struck{\ill{} plus un pt}
$H_i$ d'équations $\lambda_i = 0$ ($i \in \mathbf{Z}/n\mathbf{Z}$) --- donc,
le corps $k$ étant infini, \struck{\ill{}} ce noyau se réduit à un pt, i.e.\
est nul. Donc $N(u_*) \to k^3$ est \emph{injectif}. Donc
\[
1 \leqslant \dim N(u_*) \leqslant 3 .
\]
\note{« a un noyau » est écrit à droite de la flèche, derrière un grand
crochet ; « dans $N(u_*)$ » et « $H_i$ d'équations » sont en interligne, et la
suite de la phrase n'est reconstituée qu'au prix de plusieurs lectures
douteuses}
\marginal{Lemme ($N(u_*) \not\subset \bigcup H_i$)}
\note{note de marge, écrite en oblique, en face du passage sur les
hyperplans, avec un signe de renvoi}

Mais \struck{OPS \ill{}, \ill{}} On a $\dim N(u_*) \neq 3$, i.e.\
$N(u_*) \to k^3$ non surjectif --- \struck{il} \add{l'image} ne contient pas
p.ex.\ $(1, 0, 0)$ (car $\lambda_1 u_0 - \lambda_2 (u_1 + u_2) + \lambda_3 u_3$
impliquerait $u_0 = 0$, contraire à l'hyp.). On \uncertain{va regarder} le cas
$\dim N(u_*) = 2$. On a
\[
N(u_*) \simeq \operatorname{Im}\bigl( N(u_*) \to k^3 \bigr)
 \subset \bigl\lbrace (\lambda_1, \lambda_2, \lambda_3) \bigm|
 \lambda_1 u_0 - \lambda_2 (u_1 + u_2) + \lambda_3 u_3 = 0 \bigr\rbrace
\]
\[
= \operatorname{Ker}\bigl( (\lambda_1, \lambda_2, \lambda_3) \mapsto
 \lambda_1 u_0 - \lambda_2 (u_1 + u_2) + \lambda_3 u_3 \bigr)
\]
donc
\[
\dim N(u_*) \leqslant \dim \operatorname{Ker}
 = 3 - \operatorname{rg}(u_0, u_1 + u_2, u_3)
 = \left\lbrace
 \begin{array}{l}
 1 \text{ si } \operatorname{rg} = 2 \\
 2 \text{ si } \operatorname{rg} = 1
 \end{array}
 \right.
\]
\note{sous « $N(u_*) \simeq$ », un premier \struck{$N(u_*) \longrightarrow$} est
biffé. Dans la parenthèse « (car \ldots) », les indices et le signe devant
$\lambda_2$ sont surchargés ; le « $= 0$ » n'est pas écrit. Dans
« $3 - \operatorname{rg}$ », « rg » est écrit en surcharge de « dim »}

\page{27}
\note{p. 15 de sa pagination, numéro cerclé}

donc $\operatorname{rg} N(u_*) = 2 \Longrightarrow
\operatorname{rg}(u_i, u_{i+1} + u_{i+2}, u_{i+3}) = 1$ $\forall i$
\note{« rg $N(u_*)$ » tel qu'écrit, pour $\dim N(u_*)$}

Dém.\ a) $\Rightarrow$ b).

Un calcul \uncertain{facile} \uncertain{plus haut} montre b) $\Leftrightarrow$
c) et \uncertain{on voit de suite} que \uncertain{ceci équivaut} à d).

On considère les $(\lambda_i)$ \emph{périodiques} de période $3$, tels que
\struck{$\lambda_0 u_0 + \lambda_1 u_1 + \lambda_2$}
\[
\lambda_i u_{i-1} - \lambda_{i+1} (\underbrace{u_i + u_{i+1}}_{-u_{i-1}})
 + \underset{\textstyle \lambda_{i-1}}{\underset{\shortparallel}{\lambda_{i+2}}}
 (\underset{\textstyle u_{i-1}}{\underset{\shortparallel}{u_{i+2}}}) = 0
 \qquad \forall i
\]
i.e.
\[
(\lambda_i + \lambda_{i-1} + \lambda_{i+1})\, u_{i-1} = 0
\]
\note{le premier signe de la relation est surchargé ; dans la dernière
formule, un groupe de lettres est noirci entre $\lambda_{i-1}$ et
$\lambda_{i+1}$}

Ils forment un vectoriel de dim $2$ !

Dém.\ d) $\Rightarrow$ a)
\note{la démonstration annoncée n'est pas écrite : le reste de la page est
blanc}

\page{28}
\note{p. 16 de sa pagination, numéro cerclé}

$V$ plan vect.\ $/k$, $n$ entier $\geqslant 1$
\note{le « 1 » est écrit en surcharge d'un autre chiffre, \uncertain{un 3}}

\emph{Théorème} $u_* = (u_i)_{i \in \mathbf{Z}/n\mathbf{Z}} \in
V^{\mathbf{Z}/n\mathbf{Z}}$, les $u_i \neq 0$
\[
N(u_*) = \bigl\lbrace (\lambda_i) \in k^{\mathbf{Z}/n\mathbf{Z}} \bigm|
 \lambda_i u_{i-1} - \lambda_{i+1} (u_i + u_{i+1}) + \lambda_{i+2} u_{i+2} = 0
 \quad \forall i \in \mathbf{Z}/n\mathbf{Z} \bigr\rbrace
\]
\note{la relation est serrée et surchargée ; les indices des deuxième et
troisième $\lambda$ sont \uncertain{lus} d'après la relation $(*)$ de la
page 25}

Soit $H_i \subset k^{\mathbf{Z}/n\mathbf{Z}}$ l'hyperplan coord.\
\uncertain{$i$ème}. On a \ill{} l'un des cas \uncertain{mutuellement}
exclusifs :

(a) $\dim N(u_*) = 0$ i.e.\ $N(u_*) = 0$ [i.e., si $\sum u_i = 0$,
l'application \add{can.} $\varphi : \ill{}(V) \to (\det V)^{\mathbf{Z}/n\mathbf{Z}}$
\ill{} \emph{\uncertain{lin.}} \ill{} $u_*$]
\marginal{« Cas général »}
\note{la lettre devant « $(V)$ » est surchargée, peut-être un $X$ ;
« can. » est ajouté au-dessus de « l'application »}

(b) \struck{\ill{} $\dim N(u_*) = 1$ et}
\[
N(u_*) \not\subset \bigcup H_i \qquad
\bigl( \Longrightarrow \dim N(u_*) \in \lbrace 1, 2 \rbrace \bigr)
\]

\quad (b$_1$) $\dim N(u_*) = 1$
\marginal{« Cas spécial I »}

\quad (b$_2$) $\dim N(u_*) = 2$ \quad $\Longleftrightarrow$
$\left\lbrace
\begin{array}{l}
u_i = u_{i+3} \\
u_0 + u_1 + u_2 = 0
\end{array}
\right.$
\marginal{« Cas hyperspécial »}

(c)
$\left\lbrace
\begin{array}{l}
\exists\, i \text{ tel que } N(u_*) \subset H_i \\
N(u_*) \neq 0
\end{array}
\right.$
\marginal{« Cas spécial II »}
\note{les lettres (a), (b), (b$_1$), (b$_2$), (c) sont cerclées. Les quatre
noms de cas sont écrits en oblique dans la marge gauche, entre guillemets, et
« Cas hyperspécial » est souligné}

($\Longrightarrow \dim N(u_*) = 1$, si $\lambda = (\lambda_i) \in N(u_*)$,
$\lambda \neq 0$, \ill{} \ill{}) $\lambda_{i_0} = 0$, les \struck{\ill{}} $i$
tels que $\lambda_i = 0$ sont ceux pour lesquels $s_i = s_{i_0}$ (i.e.\
$u_{i_0} + \cdots + u_{i-1} = 0$)), \uncertain{car} $\lambda_0 = 0$,

\struck{\ill{}} Pour que \ill{} dans le cas c), il \uncertain{faut et il
suffit}
\[
\left\lbrace
\begin{array}{l}
\text{a) } \sum u_i = 0 \\
\text{b) } \forall\, 1 \leqslant i \leqslant n-1, \text{ on ait }
 u_{i-1}, u_i \parallel u_1, u_{0i}
\end{array}
\right.
\]
et si \uncertain{on normalise} ($\lambda_1 = 1$), on aura $\lambda_i$
\struck{\ill{}} défini par
\[
u_1 u_{0i} = \lambda_i\, u_i u_{i-1}
\]
\note{bas de page très serré. « $u_{0i}$ » est lu tel qu'écrit ; le sens n'en
est pas fixé sur cette page. Dans « $\forall\, 1 \leqslant i \leqslant n-1$ »
le $\forall$ est surchargé}

\marginal{$\exists\, o \in E$, \struck{\ill{}} $\Delta \in E$ \ill{} $o$), tel
que $\forall i$ tel que \struck{\ill{} $\neq 0$} $u_{i-1}$ \struck{\ill{}}
$\parallel \Delta$ et \struck{\ill{}} $u_i \parallel$ \ill{} \ldots\ $\exists\, i$,
\ill{} \ill{} \ill{}}
\note{note de marge écrite en oblique dans le coin inférieur gauche, reliée par
une double flèche $\Longleftarrow$ à « Pour que \ldots\ dans le cas c) » ;
plusieurs lignes en sont biffées et la fin, avec une numérotation 1°), 2°),
3°) \uncertain{lue}, n'est pas lisible}

\end{document}
